% =====================================================================
%  Lecture 19 -- 5G Core Architecture                     (English edition)
% =====================================================================
\documentclass[11pt,a4paper]{article}
\providecommand\ACNroot{../..}
\input{\ACNroot/templates/lecture-en}

\ACNset{number}{19}
\ACNset{title}{5G Core Network Architecture}
\ACNset{shorttitle}{5G core architecture}
\ACNset{titleother}{معماری هستهٔ نسل پنجم}
\ACNset{part}{Part IV -- 5G system and procedures (Lectures 14, 19--22)}
\ACNset{time}{6--8 hours (3 h reading, 3--5 h exercises)}
\ACNset{level}{Intermediate to advanced}
\ACNset{prereq}{Lectures 9, 12, 13, 14 and 18}
\ACNset{gen}{5}

\begin{document}
\ACNcover

% ---------------------------------------------------------------------
\section{Learning objectives}
After studying this lecture you should be able to:
\begin{objectives}
  \item \textbf{Draw} the 5G system architecture in reference-point form and \textbf{map} each NF to its EPC
        predecessor.
  \item \textbf{Explain} the 5G identities (SUPI, SUCI, 5G-GUTI, GPSI, PEI) and how SUCI protects subscriber privacy.
  \item \textbf{Describe} PDU sessions, their types and the session and service continuity (SSC) modes.
  \item \textbf{Design} UPF topologies with uplink classifiers and local anchors for edge computing, and
        \textbf{estimate} the latency and backhaul gains.
  \item \textbf{Explain} RM/CM states, registration areas and interworking with the EPC.
\end{objectives}

% ---------------------------------------------------------------------
\section{Prerequisite knowledge}
You need the EPC architecture (Lectures~9, 12, 13), the 5G overview (Lecture~14) and the service-based mechanics of
Lecture~18 (NRF, SBI). This lecture describes \emph{what} the 5G core consists of; Lecture~20 describes the procedures that
use it, Lecture~21 slicing and Lecture~22 QoS.

\paragraph{Key terms} The vocabulary of this lecture with its Persian equivalents:
\begin{keyterms}
SUPI / SUCI & \fa{شناسهٔ دائمی / پنهان‌شدهٔ مشترک} & Permanent identity / its encrypted form sent over the air. \\
5G-GUTI & \fa{شناسهٔ موقت سراسری} & Temporary identity allocated by the AMF. \\
PDU session & \fa{نشست PDU} & Association between UE and a data network through the 5GC. \\
DNN & \fa{نام شبکهٔ داده} & Name of a data network (the 5G equivalent of an APN). \\
SSC mode & \fa{حالت پیوستگی نشست و خدمت} & Rules for keeping or changing the session anchor during mobility. \\
PSA & \fa{لنگر نشست PDU} & UPF terminating N6 towards the data network. \\
UL CL / BP & \fa{طبقه‌بند فراسو / نقطهٔ انشعاب} & UPF functions that split traffic between anchors. \\
RM / CM & \fa{مدیریت ثبت / اتصال} & Registration and connection management states. \\
N26 & \fa{واسط N26} & AMF--MME interface for EPC--5GC interworking. \\
N3IWF & \fa{تابع همکاری غیر 3GPP} & Gateway for untrusted non-3GPP access (e.g.\ Wi-Fi). \\
\end{keyterms}

% ---------------------------------------------------------------------
\section{Introduction and motivation}
The 5G core (5GC) is the first mobile core designed from the outset for cloud deployment. It keeps the proven concepts of
the EPC -- a mobility anchor in the control plane, a session anchor in the user plane, a subscriber database, policy
control -- but reorganises them: access and mobility are separated from session management, the user plane is fully
controlled through PFCP, subscription data are split between front-ends and a repository, and every control-plane function
is a service producer (Lecture~18) \cite{ts23501}. The new structure supports features the EPC could not offer natively:
access-agnostic registration (3GPP and non-3GPP), multiple user-plane anchors per session for edge computing, Ethernet and
unstructured sessions, network slicing and a flow-based QoS model.

% ---------------------------------------------------------------------
\section{Fundamental concepts}

\subsection{Network functions and reference points}
\Cref{fig:ref} shows the 5G system in reference-point representation. Each reference point is realised on the SBI by the
services of the producer (e.g.\ N11 by Nsmf services, N8 by Nudm services), except N1, N2, N3, N4, N6 and N9.
\begin{table}[htbp]
\centering\small
\caption{5GC network functions and their EPC counterparts.}\label{tab:map}
\rowcolors{2}{acnmist}{white}
\begin{tabularx}{\linewidth}{@{}l Y l@{}}
\toprule
\rowcolor{white}5GC & Main responsibilities & EPC counterpart \\
\midrule
AMF  & Registration, connection, reachability, mobility; NAS termination and security; SMS & MME (mobility part) \\
SMF  & PDU sessions, IP allocation, UPF selection and control, charging trigger & MME (session part), SGW-C, PGW-C \\
UPF  & Packet routing, forwarding, QoS enforcement, usage reporting, anchoring & SGW-U, PGW-U \\
AUSF & Authentication server                            & HSS (AuC part) \\
UDM / UDR & Credentials processing, subscription data / repository & HSS \\
PCF  & Policy for access, mobility, sessions and UE     & PCRF \\
NSSF & Slice selection                                  & -- \\
NEF  & Exposure of capabilities to external applications & SCEF \\
NRF, SCP, SEPP & Discovery, indirect communication, roaming security & DNS, DRA, DEA \\
\bottomrule
\end{tabularx}
\end{table}

\begin{acnfigure}
  \resizebox{0.95\linewidth}{!}{\acnfig{lecture-19/ref-points}}
  \caption{5G system architecture in reference-point representation (non-roaming, selected NFs) \cite{ts23501}.}
  \label{fig:ref}
\end{acnfigure}

\subsection{Identities}
\begin{itemize}
  \item \textbf{SUPI}: the permanent subscription identifier, usually an IMSI (or a network access identifier).
  \item \textbf{SUCI}: the SUPI concealed with the home network's public key; only the routing part is in clear
        (\cref{fig:ids}). The UDM's de-concealment function recovers the SUPI \cite{ts33501}.
  \item \textbf{5G-GUTI}: temporary identity = GUAMI (PLMN, AMF region, AMF set, AMF pointer) + 5G-TMSI; the 48-bit
        5G-S-TMSI is used in paging and RRC.
  \item \textbf{PEI}: permanent equipment identifier (IMEI); \textbf{GPSI}: generic public identifier (MSISDN or external ID)
        used towards external applications.
\end{itemize}
Unlike LTE, a 5G UE never sends its permanent identity in clear over the radio interface: even the first registration uses a
SUCI, which closes the ``IMSI catcher'' weakness of earlier generations.

\begin{acnfigure}
  \resizebox{\linewidth}{!}{\acnfig{lecture-19/identifiers}}
  \caption{Structure of the SUCI and of the 5G-GUTI.}
  \label{fig:ids}
\end{acnfigure}

\subsection{PDU sessions}
A PDU session connects the UE to a data network identified by a DNN and a slice (S-NSSAI). Session types are IPv4, IPv6,
IPv4v6, \emph{Ethernet} (the UPF switches Ethernet frames, e.g.\ for industrial LANs) and \emph{Unstructured} (opaque data).
Within a session, traffic is carried in QoS flows (Lecture~22); a UE may have several sessions to different DNNs or slices at
the same time.

\subsection{Session and service continuity}
The \emph{SSC mode} of a session defines what happens to its anchor when the UE moves \cite{ts23501}:
\begin{itemize}
  \item \textbf{SSC mode 1}: the PSA (and the IP address) is kept for the life of the session -- like the EPC's P-GW.
  \item \textbf{SSC mode 2}: the network may release the session and ask the UE to establish a new one with a closer anchor
        (``break before make'').
  \item \textbf{SSC mode 3}: the network establishes a new session with a new anchor before releasing the old one (``make
        before break''); the application can migrate gracefully.
\end{itemize}

% ---------------------------------------------------------------------
\section{Detailed technical explanation}

\subsection{The AMF}
The AMF terminates N1 (NAS) and N2 (NGAP) \cite{ts24501,ts38413}. It handles registration, reachability and paging,
connection management, mobility and access authentication with the AUSF. Unlike the MME it does not process session
signalling: session NAS messages are relayed transparently to the SMF. One AMF serves a UE across 3GPP and non-3GPP
access, so registration is access-agnostic. AMFs are organised in \emph{AMF sets}: any AMF of a set can take over a UE's
context from shared storage.

\subsection{The SMF and UPF}
The SMF manages PDU sessions: it selects UPFs, allocates IP addresses, installs PFCP rules (PDR, FAR, QER, URR, BAR -- Lecture~13)
over N4, interacts with the PCF for policy and with the CHF for charging. The UPF is the only user-plane function: an
\emph{intermediate UPF} (I-UPF) may connect the RAN, an \emph{uplink classifier} (UL CL) or \emph{branching point} (for IPv6
multi-homing) may divert part of the traffic, and one or more \emph{PDU session anchors} (PSAs) terminate N6 (\cref{fig:upf}).

\begin{acnfigure}
  \resizebox{0.95\linewidth}{!}{\acnfig{lecture-19/upf-topology}}
  \caption{A PDU session with an uplink classifier and two anchors: local edge breakout and central Internet access.}
  \label{fig:upf}
\end{acnfigure}

\subsection{Edge computing}
To serve low-latency applications, the SMF can insert a UL CL and a local PSA at an edge site, so that traffic for edge
application addresses is routed locally while the rest goes to the central anchor. Application functions can influence this
routing through the NEF or PCF, and TS~23.548 adds discovery of edge application servers \cite{ts23548}. The result is a single
PDU session with both local and central breakout -- impossible in the EPC without special arrangements.

\subsection{States and registration areas}
RM states: RM-DEREGISTERED and RM-REGISTERED; CM states: CM-IDLE and CM-CONNECTED. A UE in CM-CONNECTED may additionally be in
RRC\_INACTIVE: the N2 connection and context stay in place while the RAN pages the UE within a RAN notification area, which
saves signalling for frequent small transmissions. The AMF allocates a \emph{registration area} (a TAI list) and performs
mobility and periodic registration updates, as the MME did with TAUs.

\subsection{Interworking with the EPC}
During the long coexistence of 4G and 5G, a UE moves between LTE (EPC) and NR (5GC). With \emph{N26} the MME and AMF exchange
contexts directly and IP sessions survive the transition; this requires combined nodes: SMF+PGW-C, UPF+PGW-U, HSS+UDM and
PCF+PCRF serve both systems (\cref{fig:iw}). Without N26 the UE re-registers and the session continuity depends on the
combined anchor. Voice is supported by EPS fallback (a 5G voice call is handed over to LTE for VoLTE) where VoNR is not
available.

\begin{acnfigure}
  \resizebox{\linewidth}{!}{\acnfig{lecture-19/interworking}}
  \caption{EPC--5GC interworking with N26 and combined nodes.}
  \label{fig:iw}
\end{acnfigure}

\subsection{Non-3GPP access}
Untrusted non-3GPP access (e.g.\ public Wi-Fi) connects through an N3IWF, which terminates an IPsec tunnel from the UE and
presents N2/N3 towards the core; trusted non-3GPP access uses a TNGF. In both cases the UE registers with the same AMF and
uses the same NAS and PDU session procedures -- the access-agnostic design of the 5GC.

% ---------------------------------------------------------------------
\section{Mathematical models}

\subsection{SUCI size}\label{sec:suci}
\paragraph{Model} With ECIES, the scheme output consists of the ephemeral public key ($k$ bytes), the encrypted MSIN (the same
length $m$ as the plaintext, BCD encoded) and a MAC tag ($t$ bytes):
\begin{equation}\label{eq:suci}
  L_{\mathrm{out}} = k + m + t .
\end{equation}
\paragraph{Example} A 10-digit MSIN ($m = 5$ bytes): profile A (X25519, $k = 32$, $t = 8$) gives 45 bytes; profile B (P-256
compressed point, $k = 33$) gives 46 bytes \cite{ts33501}. The overhead is small compared with the privacy gained, and the
computation is done once per registration attempt that needs a permanent identity.

\subsection{Latency and UPF placement}\label{sec:rtt}
\paragraph{Model} Round-trip latency between UE and an application server behind a PSA at fibre distance $D$ is
\begin{equation}\label{eq:rtt}
  \mathrm{RTT} \approx \mathrm{RTT}_{\mathrm{radio}} + 2\,\frac{D}{v} + T_{\mathrm{proc}},
\end{equation}
with $v \approx 2\times10^8$\,m/s in fibre (about 5\,\textmu s per km).
\paragraph{Example} $D = 500$\,km (central data centre) adds 5\,ms; $D = 20$\,km (edge site) adds 0.2\,ms. With a 4\,ms radio RTT
and 1\,ms processing, the RTT falls from 10\,ms to 5.2\,ms -- the edge halves it.
\paragraph{Interpretation} Propagation is only one term; for URLLC targets of a few milliseconds every component -- radio
numerology, scheduling, transport and processing -- must be engineered together.

\subsection{Backhaul offload}\label{sec:offload}
\paragraph{Model} If a fraction $\alpha$ of a site's traffic $T$ is served locally (UL CL to a local PSA and edge cache or
application), the traffic carried towards the central core is
\begin{equation}\label{eq:offload}
  T_{\mathrm{bh}} = (1-\alpha)\,T .
\end{equation}
\paragraph{Example} With 60\,\% of video served from an edge cache, a site with 10\,Gbit/s busy-hour traffic needs 4\,Gbit/s of
backhaul towards the central core instead of 10.

\subsection{Interruption in SSC modes}\label{sec:ssc}
\paragraph{Model} When the anchor changes, the service interruption is approximately
\begin{equation}\label{eq:ssc}
  T_{\mathrm{int}} \approx
  \begin{cases}
    0 & \text{SSC 1 (no anchor change) and SSC 3 (make before break)},\\
    T_{\mathrm{rel}} + T_{\mathrm{est}} + T_{\mathrm{app}} & \text{SSC 2 (break before make)},
  \end{cases}
\end{equation}
where $T_{\mathrm{rel}}$ and $T_{\mathrm{est}}$ are the release and establishment times and $T_{\mathrm{app}}$ the time for the
application to reconnect with the new address.
\paragraph{Example} $T_{\mathrm{rel}} = 50$\,ms, $T_{\mathrm{est}} = 150$\,ms, $T_{\mathrm{app}} = 300$\,ms: half a second without
service in SSC~2 -- acceptable for a background download, not for a video call, which should use SSC~1 or~3.

% ---------------------------------------------------------------------
\section{Network architecture and diagrams}
A national 5GC typically has a few central sites hosting AMF sets, SMFs, UDM/UDR, PCF, NRF and the central UPFs, and many edge
sites with UPFs (and possibly edge application platforms) close to the RAN. The SMF chooses, per session, which UPFs to chain:
for an Internet session a central PSA, for an industrial Ethernet session a local PSA on the factory premises, for a session
with edge services a UL CL with both (\cref{fig:upf}). Slices (Lecture~21) may use dedicated SMFs and UPFs.

% ---------------------------------------------------------------------
\section{Step-by-step procedures}

\begin{procedurebox}{SUCI calculation and de-concealment}\label{proc:suci}
\begin{steps}
  \item The USIM (or ME) holds the home network public key and its identifier, and the protection scheme to use.
  \item For each registration that needs a permanent identity, the UE generates an ephemeral key pair, derives a shared secret
        with the home network public key and encrypts the MSIN; it computes the MAC tag.
  \item The UE sends the SUCI (routing part in clear plus scheme output) in the Registration Request.
  \item The AMF forwards it via the AUSF to the home UDM, whose SIDF uses the private key to recover the SUPI.
\end{steps}
\end{procedurebox}

\begin{procedurebox}{Adding a local anchor for edge traffic}\label{proc:ulcl}
\begin{steps}
  \item An edge application is deployed and its address range is configured in (or reported through the NEF/PCF to) the SMF.
  \item When the UE is in the edge site's area, the SMF selects a UPF there as UL CL and a local PSA, and installs PFCP rules: traffic
        matching the edge addresses is forwarded to the local PSA, the rest over N9 to the central PSA.
  \item The SMF updates the RAN's N3 tunnel to point to the UL CL (via the AMF and N2).
  \item When the UE leaves the area, the SMF removes the UL CL and local PSA; the session continues through the central anchor.
\end{steps}
\end{procedurebox}

% ---------------------------------------------------------------------
\section{Worked examples}

\begin{examplebox}{Mapping EPC to 5GC}\label{ex:map}
\textbf{Problem.} Which 5GC functions take over (a) the MME's tracking-area list management, (b) the MME's bearer set-up via S11,
(c) the P-GW's IP address allocation, (d) the PCRF's Gx rules?

\textbf{Solution.} (a) AMF (registration areas); (b) SMF (PDU session management, N4 to the UPF); (c) SMF (allocation) with the
UPF as anchor; (d) PCF over N7 (Npcf\_SMPolicyControl) to the SMF.
\end{examplebox}

\begin{examplebox}{Edge latency budget}\label{ex:edge}
\textbf{Problem.} An AR application needs an RTT below 10\,ms. Radio RTT is 4\,ms, processing 1\,ms. Up to which fibre distance can the
application server be placed?

\textbf{Solution.} $4 + 1 + 2D/v \le 10 \Rightarrow D \le 5\ \text{ms}\times v/2 = 500$\,km of fibre path. Fibre routes are typically 1.5 times
the straight-line distance, so about 330\,km in a straight line -- a regional, not necessarily a cell-site, edge.
\end{examplebox}

\begin{examplebox}{Choosing an SSC mode}\label{ex:ssc}
\textbf{Problem.} Choose SSC modes for (a) IMS voice, (b) a vehicle downloading map updates, (c) a cloud game that can migrate between edge
servers.

\textbf{Solution.} (a) SSC~1 -- the IMS registration is bound to the IP address; (b) SSC~2 -- a short interruption is tolerable and a local
anchor reduces backhaul; (c) SSC~3 -- the application opens a connection via the new anchor before the old one is released.
\end{examplebox}

% ---------------------------------------------------------------------
\section{Technology comparison}
\begin{table}[htbp]
\centering\small
\caption{EPC and 5GC architecture compared.}\label{tab:cmp}
\rowcolors{2}{acnmist}{white}
\begin{tabularx}{\linewidth}{@{}l Y Y@{}}
\toprule
\rowcolor{white} & EPC & 5GC \\
\midrule
Mobility vs session control & Both in MME          & AMF and SMF separated \\
User plane control & Integrated (CUPS optional)    & Always via PFCP (N4) \\
Session types      & IP (non-IP via SCEF)          & IP, Ethernet, Unstructured \\
Anchors per session & One P-GW                     & One or more PSAs, UL CL/BP \\
Continuity         & Anchor fixed                  & SSC modes 1, 2, 3 \\
Identity privacy   & IMSI may be sent in clear     & SUCI (always concealed) \\
Access             & 3GPP; non-3GPP via ePDG       & Access-agnostic AMF (N3IWF, TNGF) \\
Signalling         & Diameter, GTPv2-C             & SBI (HTTP/2), NAS, NGAP, PFCP \\
\bottomrule
\end{tabularx}
\end{table}

\paragraph{Continuity of ideas} Many EPC concepts survive with new names: TAI lists become registration areas, APNs become DNNs, bearers
become QoS flows, the HSS becomes UDM/UDR/AUSF. Understanding the EPC therefore remains the best preparation for the 5GC.

% ---------------------------------------------------------------------
\section{Key takeaways}
\begin{takeawaybox}
\begin{itemize}
  \item 5GC separates AMF (access, mobility) from SMF (sessions) and controls UPFs only through PFCP.
  \item SUPI is never sent in clear; SUCI = routing info + ECIES output ($k + m + t$ bytes).
  \item PDU sessions (IP, Ethernet, Unstructured) can have several anchors; UL CL and BP enable local breakout.
  \item SSC modes 1/2/3 trade anchor stability against optimal routing; SSC~2 interrupts service, SSC~3 does not.
  \item Edge UPFs cut RTT by $2D/v$ and backhaul by a factor $(1-\alpha)$.
  \item Interworking with the EPC relies on N26 and combined nodes (SMF+PGW-C, UPF+PGW-U, HSS+UDM, PCF+PCRF).
\end{itemize}
\end{takeawaybox}

% ---------------------------------------------------------------------
\section{Exercises}

\begin{exercisebox}{1}{Who does it?}\label{exr:who}
Which NF performs: (a) paging; (b) IP address allocation; (c) slice selection at registration; (d) SUCI de-concealment;
(e) exposure of UE location to an authorised external application?
\end{exercisebox}

\begin{exercisebox}{1}{Ethernet sessions}\label{exr:eth}
Give an application of Ethernet PDU sessions and explain why IP sessions would be inconvenient for it.
\end{exercisebox}

\begin{exercisebox}{2}{SUCI size}\label{exr:suci}
Compute the scheme-output length for a 9-digit MSIN under profiles A and B. How many bytes does concealment add compared with the
null scheme, which sends the MSIN in clear?
\end{exercisebox}

\begin{exercisebox}{2}{Edge offload}\label{exr:offload}
A city has 50 edge sites, each with 8\,Gbit/s busy-hour traffic. What edge-served fraction is needed to keep total backhaul to the
central core below 200\,Gbit/s?
\end{exercisebox}

\begin{exercisebox}{3}{Latency design}\label{exr:latency}
A factory needs RTT below 5\,ms between a robot controller and the UE. Radio RTT is 2\,ms with mini-slots and processing 0.5\,ms. Where must
the PSA and the controller be, and which session type and SSC mode would you choose? What if the radio RTT were 4\,ms?
\end{exercisebox}

\begin{exercisebox}{3}{Interworking without N26}\label{exr:n26}
Explain what happens to an ongoing IP session when a UE moves from NR to LTE coverage (a) with N26 and combined anchors, (b) without
N26 but with combined anchors, (c) without combined anchors. Which case keeps the IP address?
\end{exercisebox}

\subsection*{Solutions}
\addcontentsline{toc}{subsection}{Solutions}

\begin{solutionbox}{\ref{exr:who}}
(a) AMF (and the RAN for RRC\_INACTIVE); (b) SMF; (c) AMF with the NSSF; (d) UDM (SIDF); (e) NEF, obtaining the information from the AMF/UDM.
\end{solutionbox}

\begin{solutionbox}{\ref{exr:eth}}
Industrial LANs (e.g.\ time-sensitive networking between controllers and machines) and enterprise Layer-2 VPNs. The devices use Ethernet
protocols and MAC addressing, sometimes non-IP; carrying them over IP sessions would need tunnelling or gateways, adding latency and
complexity, whereas an Ethernet session lets the UPF switch frames directly.
\end{solutionbox}

\begin{solutionbox}{\ref{exr:suci}}
9 digits in BCD occupy 5 bytes (last half-byte padded). Profile A: $32 + 5 + 8 = 45$ bytes; profile B: $33 + 5 + 8 = 46$ bytes. The null scheme
carries the 5-byte MSIN, so concealment adds 40 or 41 bytes.
\end{solutionbox}

\begin{solutionbox}{\ref{exr:offload}}
Total traffic $50\times8 = 400$\,Gbit/s; $(1-\alpha)\times400 \le 200 \Rightarrow \alpha \ge 0.5$: half of the traffic must be served at the edge.
\end{solutionbox}

\begin{solutionbox}{\ref{exr:latency}}
$2 + 0.5 + 2D/v \le 5 \Rightarrow D \le 250$\,km: in principle a regional site, but transport hops, queueing and jitter consume budget, so the PSA
and controller are placed on the factory premises (local UPF). An Ethernet session (if the control protocol is Layer-2) with SSC~1 (fixed
local anchor). With 4\,ms radio RTT only 1\,ms remains: $D \le 100$\,km, and in practice the same on-premises design, with little margin.
\end{solutionbox}

\begin{solutionbox}{\ref{exr:n26}}
(a) The AMF passes the context to the MME over N26; the session is moved to the EPC through the same SMF+PGW-C/UPF+PGW-U: IP address
kept, short interruption. (b) The UE attaches/registers in the EPC itself with ``handover'' indications; the combined anchor can restore the
PDN connection with the same address, but with a longer interruption. (c) A new PDN connection with a new IP address: the application must
reconnect. Only (a) and (b) keep the IP address.
\end{solutionbox}

% ---------------------------------------------------------------------
\section{Further reading}
TS~23.501 \cite{ts23501} specifies the architecture and TS~23.502 \cite{ts23502} the procedures; TS~23.503 \cite{ts23503} policy;
TS~23.548 \cite{ts23548} edge computing; TS~33.501 \cite{ts33501} security including SUCI; TS~24.501 \cite{ts24501} NAS; TS~38.413
\cite{ts38413} NGAP. Rommer et al.\ \cite{rommer2020core} give a thorough textbook account of the 5G core.

\printbibliography[heading=subbibliography,title={\ACNstrRefs}]
\end{document}
